\section*{अध्याय \ref{ch:g8:negprod} --- ऋण संख्याओं का गुणा}

\begin{solution}{exo:g8:negprod:1}
$(-8) \times (+7) = -56$; \quad
$(-9) \times (-6) = +54$; \quad
$(+12) \times (-0.5) = -6$; \quad
$(-1) \times (-1) \times (-1) = -1$ (तीन ऋण गुणक: विषम)।
\end{solution}

\begin{solution}{exo:g8:negprod:2}
$(-63) \div (+9) = -7$; \quad
$(-8) \div (-16) = +0.5$; \quad
$\dfrac{45}{-5} = -9$; \quad
$\dfrac{-7.2}{-8} = +0.9$.
\end{solution}

\begin{solution}{exo:g8:negprod:3}
तीन ऋण गुणक: \emph{ऋण}। \quad
$(-1)^{10}$: दस ऋण गुणक, सम: \emph{धन} (वह $1$ के बराबर है)।
\quad
$(-2)^7$: सात गुणक, विषम: \emph{ऋण}। \quad
$(-4) \times 0 \times (-6) = 0$: न धन, न ऋण।
\end{solution}

\begin{solution}{exo:g8:negprod:4}
$(-2)^4 = +16$; \quad $-2^4 = -16$ (ऋण चिह्न वर्ग नहीं हुआ); \quad
$(-3)^3 = -27$; \quad $(-1)^{2026} = +1$ (सम घातांक)।
\end{solution}

\begin{solution}{exo:g8:negprod:5}
$7 + 2 \times (-6) = 7 + (-12) = -5$.

$(-4) \times (5 - 9) = (-4) \times (-4) = +16$.

$18 \div (-3) - (-2) = -6 + 2 = -4$.
\end{solution}

\begin{solution}{exo:g8:negprod:6}
$\dfrac{(-5) \times (+8)}{-4} = \dfrac{-40}{-4} = +10$.

$\dfrac{(-3) \times (-6)}{(-2) \times (+9)} = \dfrac{+18}{-18} = -1$.
\end{solution}

\begin{solution}{exo:g8:negprod:7}
$x = -3$ पर: \ $5x = -15$; \quad $x^2 = (-3)^2 = 9$; \quad
$-x^2 = -9$; \quad $2x^2 + 4x = 18 - 12 = 6$; \quad
$(2x)^2 = (-6)^2 = 36$।
\end{solution}

\begin{solution}{exo:g8:negprod:8}
$(-6) \times (-7) = 42$; \quad
$10 \div (-4) = -2.5$; \quad
$(-5) \times \frac15 = -1$, इसलिए छूटी हुई संख्या
$+\frac15 = 0.2$ है।
\end{solution}

\begin{solution}{exo:g8:negprod:9}
\emph{1. ग़लत}: $0.5 \times 0.5 = 0.25$, $0.5 + 0.5 = 1$ से छोटा है
(या: $2 \times 3 = 6 > 5$, पर दावे को हमेशा सही होना चाहिए)।

\emph{2. सही}: वर्ग एक ही चिह्न वाली दो संख्याओं का \omterm{def:g2:mult:def}{गुणनफल} है,
इसलिए वह धन या \omterm{def:g1:counting:zero}{शून्य} होता है।

\emph{3. ग़लत}: $(-2) \times (-3) \times 4 = +24$ दो ऋण गुणकों के
साथ भी धन है।
\end{solution}

\begin{solution}{exo:g8:negprod:10}
मान लो सही उत्तरों की संख्या $r$ है; तब $20 - r$ ग़लत हैं, और
\[
5r - 3(20 - r) = 36
\ \Longrightarrow\
5r - 60 + 3r = 36
\ \Longrightarrow\
8r = 96
\ \Longrightarrow\
r = 12 .
\]
बारह उत्तर सही थे (और आठ ग़लत: जाँच,
$60 - 24 = 36$)।
\end{solution}

\begin{solution}{exo:g8:negprod:11}
\omterm{def:g1:counting:zero}{शून्य} गुना कुछ भी \omterm{def:g1:counting:zero}{शून्य} होता है, इसलिए $(-1) \times 0 = 0$।
$0 = 1 + (-1)$ लिखकर वितरण नियम लगाने पर:
\[
0 = (-1) \times \bigl(1 + (-1)\bigr)
= (-1) \times 1 + (-1) \times (-1)
= -1 + (-1) \times (-1).
\]
$(-1) \times (-1)$ को $-1$ में जोड़ने पर $0$ मिलता है: इसलिए उसे
$-1$ की \omterm{def:g7:negatives:def}{विपरीत संख्या}, यानी $+1$, होना ही पड़ेगा। वितरण नियम कोई
दूसरा रास्ता छोड़ता ही नहीं।
\end{solution}

\begin{solution}{pb:g8:negprod:1}
\textbf{1.} चूँकि $0 = 0 + 0$ है, वितरण नियम
(\cref{thm:g7:priorities:distributivity}) से
\[
0 \times a = (0 + 0) \times a = 0 \times a + 0 \times a .
\]
यानी $0 \times a$ को अपने आप में जोड़ने से कुछ नहीं बदलता --- और बिना
कुछ बदले जोड़ी जा सकने वाली अकेली संख्या $0$ है। (सीधे कहें: दोनों
तरफ़ से $0 \times a$ घटा दो।) इसलिए $0 \times a = 0$।

\textbf{2.} फैलाकर, फिर $1 + (-1) = 0$ और सवाल 1 का इस्तेमाल करके:
\[
a + (-1) \times a
= 1 \times a + (-1) \times a
= \bigl(1 + (-1)\bigr) \times a
= 0 \times a = 0 .
\]
यानी $(-1) \times a$ वह संख्या है जो $a$ में जोड़ने पर $0$ देती है:
और वह ठीक $a$ की \omterm{def:g7:negatives:def}{विपरीत संख्या} है, इसलिए $(-1) \times a = -a$।
\Cref{exo:g8:negprod:11} इसकी ख़ास स्थिति $a = -1$ है: वहाँ
$(-1) \times (-1) = -(-1) = +1$।

\textbf{3.} $-a = (-1) \times a$ लिखो (सवाल 2) और गुणकों को नए सिरे
से जमाओ:
\[
(-a) \times b
= \bigl((-1) \times a\bigr) \times b
= (-1) \times (a \times b)
= -(a \times b) .
\]

\textbf{4.} सवाल 3 को दो बार लगाओ (एक-एक गुणक पर):
\[
(-a) \times (-b)
= -\bigl(a \times (-b)\bigr)
= -\bigl(-(a \times b)\bigr)
= a \times b ,
\]
क्योंकि किसी संख्या की विपरीत की विपरीत वही संख्या होती है।

\textbf{5.} $a = +d$ या $a = -d$ और $b = +e$ या $b = -e$ लिखो, जहाँ
$d$ और $e$ \omterm{def:g1:counting:zero}{शून्य} से दूरियाँ हैं। चारों मामले:
\[
(+d) \times (+e) = de, \quad
(+d) \times (-e) = -de, \quad
(-d) \times (+e) = -de, \quad
(-d) \times (-e) = +de,
\]
जिनमें बीच के दो सवाल 3 से और आख़िरी सवाल 4 से मिलते हैं। हर मामले
में उत्तर की दूरी $d \times e$ है, यानी दूरियों का \omterm{def:g2:mult:def}{गुणनफल}, और चिह्न
वही है जो \cref{thm:g8:negprod:rules} की सारणी बताती है: एक जैसे
चिह्न पर धन, उलटे चिह्न पर ऋण।

\textbf{6.} $(-1)^2 = +1$, $(-1)^3 = -1$, $(-1)^4 = +1$,
$(-1)^5 = -1$। आम तौर पर $(-1)^n$, $1$ दूरी वाले $n$ ऋण गुणकों का
\omterm{def:g2:mult:def}{गुणनफल} है: \cref{prop:g8:negprod:many} से वह $n$ के सम होने पर $+1$
और विषम होने पर $-1$ होता है।

\textbf{7.} ऋण गुणक दस हैं --- यानी \omterm{def:g3:numbers:evenodd}{सम संख्या} --- इसलिए \omterm{def:g2:mult:def}{गुणनफल}
\emph{धन} है (\cref{prop:g8:negprod:many})। उसकी दूरी
$1 \times 2 \times 3 \times \dots \times 10$ है।

\textbf{8.} सवाल 4 से $(-a)^2 = (-a) \times (-a) = a \times a =
a^2$। फिर सवाल 3 से,
\[
(-a)^3 = (-a)^2 \times (-a) = a^2 \times (-a) = -(a^2 \times a)
= -a^3 .
\]
सम घातांकों पर ऋण चिह्न ग़ायब हो जाता है और विषम घातांकों पर बचा रहता
है, ठीक \cref{ex:g8:negprod:powers} के मुताबिक़।

\textbf{9.} अगर $a = 0$ हो तो $a^2 = 0$। नहीं तो $a$ और $a$ के
चिह्न \emph{एक जैसे} हैं, इसलिए $a^2 = a \times a$ धन होता है
(\cref{thm:g8:negprod:rules})। इसलिए वर्ग कभी ऋण नहीं होता; और $-4$
ऋण है, इसलिए कोई भी \omterm{def:g7:negatives:def}{सापेक्ष संख्या} $x$, $x^2 = -4$ नहीं मानती।

\textbf{10.} ज़ोए सही है। घातांक $2026$ और $2027$ की सम-विषमता उलटी
है, इसलिए $(-1)^{2026} = +1$ और $(-1)^{2027} = -1$, जिनका योग $0$
है। आम तौर पर दो लगातार घातांकों में एक सम और एक विषम होता है,
इसलिए दोनों घातें किसी न किसी क्रम में $+1$ और $-1$ होती हैं:
\omterm{def:g7:negatives:def}{विपरीत संख्याएँ}, जिनका योग हमेशा $0$ होता है।

\textbf{11.} $A_3 = 1 - 2 + 3 = 2$, फिर $A_4 = A_3 - 4 = -2$,
$A_5 = A_4 + 5 = 3$, $A_6 = A_5 - 6 = -3$, $A_7 = A_6 + 7 = 4$।
अटकल: सम $n$ के लिए $A_n = -\frac{n}{2}$, और विषम $n$ के लिए
$A_n = \frac{n+1}{2}$।

\textbf{12.} सम $n$ के लिए पद पूरी तरह जोड़ों में बँट जाते हैं:
\[
A_n = (1 - 2) + (3 - 4) + \dots + \bigl((n-1) - n\bigr),
\]
और प्रत्येक जोड़ा $-1$ के बराबर है। जोड़े $\frac{n}{2}$ हैं (हर जोड़े
में दो पद, और कुल $n$ पद), इसलिए
$A_n = \frac{n}{2} \times (-1) = -\frac{n}{2}$।

\textbf{13.} विषम $n$ के लिए योग $A_n$, पहले $n - 1$ पदों का योग
$A_{n-1}$ है और उसमें आख़िरी पद जुड़ा हुआ, जो $+n$ है क्योंकि $n$
विषम है। और $n - 1$ सम है, इसलिए सवाल 12 से
\[
A_n = A_{n-1} + n = -\frac{n-1}{2} + n
= \frac{-(n-1) + 2n}{2} = \frac{n+1}{2} .
\]

\textbf{14.} $A_{100} = -\frac{100}{2} = -50$, \quad
$A_{101} = \frac{102}{2} = 51$, \quad
$A_{2026} = -\frac{2026}{2} = -1013$।

\textbf{15.} इसमें हैरानी की कोई बात नहीं: $A_{100}$ से $A_{101}$
पर जाने में ठीक एक पद जुड़ता है, यानी $+101$, इसलिए \omterm{ex:g1:subtraction:difference}{अंतर} $101$ ही
होना था। इसी कारण $A_{2027} - A_{2026} = +2027$ --- और $A_{2027}$
निकालने की कोई ज़रूरत ही नहीं।

\end{solution}
